% !TEX encoding = UTF-8
\chapter{Ứng dụng Zero Trust bảo mật hệ thống Microservices}

%==============================================================================
% 2.1 BÀI TOÁN BẢO MẬT ĐẶC THÙ
%==============================================================================
\section{Bài toán bảo mật đặc thù của hệ thống Microservices}

Kiến trúc Microservices là một dạng cụ thể của kiến trúc phân tán, trong đó ứng dụng được tổ chức thành tập hợp các dịch vụ nhỏ, độc lập, mỗi dịch vụ chạy trong tiến trình riêng và giao tiếp qua các cơ chế nhẹ như HTTP hoặc gRPC \parencite{lewis_fowler_2014}. Kiến trúc này mang lại khả năng triển khai độc lập và mở rộng linh hoạt, nhưng đồng thời làm thay đổi cơ bản bề mặt tấn công so với kiến trúc nguyên khối.

Trong phạm vi đồ án, Kubernetes được chọn làm nền tảng điều phối đại diện cho lớp container orchestration. Bài toán bảo mật vì vậy cần được phân tích ở hai tầng độc lập: rủi ro nội tại của kiến trúc Microservices và rủi ro cộng thêm từ môi trường container hóa.

%------------------------------------------------------------------------------
\subsection{Rủi ro từ kiến trúc Microservices}
%------------------------------------------------------------------------------

Kiến trúc Microservices sinh ra các thách thức bảo mật mang tính hệ thống, độc lập với môi trường triển khai:

\begin{itemize}[itemsep=6pt, parsep=0pt]
    \item \textbf{Gia tăng lưu lượng nội bộ:}
    Trong ứng dụng nguyên khối, các thành phần giao tiếp qua lời gọi hàm trong bộ nhớ. Trong Microservices, mọi tương tác đều đi qua mạng. Nếu hệ thống duy trì cơ chế tin cậy mặc định cho lưu lượng nội bộ, một dịch vụ bị xâm phạm có thể trở thành điểm xuất phát cho hành vi dò quét và di chuyển ngang đến toàn bộ hệ thống.

    \item \textbf{Phân mảnh thông tin xác thực:}
    Số lượng dịch vụ tỷ lệ thuận với số lượng thông tin xác thực cần quản lý -- khóa API, token, chứng chỉ TLS, khóa mã hóa và thông tin kết nối dịch vụ. Việc lưu trữ tĩnh các thông tin này trong mã nguồn, tệp cấu hình hoặc biến môi trường tạo ra nhiều điểm rò rỉ tiềm năng và làm tăng phạm vi ảnh hưởng khi một thông tin xác thực bị đánh cắp.
\end{itemize}

%------------------------------------------------------------------------------
\subsection{Rủi ro từ môi trường container và nền tảng điều phối}
%------------------------------------------------------------------------------

Khi workload được đóng gói bằng container và vận hành qua Kubernetes, các thành phần hạ tầng này tạo thêm các điểm yếu ở tầng hệ thống:

\begin{itemize}[itemsep=6pt, parsep=0pt]
    \item \textbf{Chia sẻ nhân hệ điều hành:}
    Khác với máy ảo, các container trên cùng một Node chia sẻ chung nhân Linux. Ranh giới namespace chỉ mang tính cách ly logic. Lỗ hổng tại tầng container runtime có thể bị khai thác để thoát khỏi vùng chứa và leo thang đặc quyền lên máy chủ vật lý.

    \item \textbf{Vòng đời ngắn của workload:}
    Kubernetes liên tục tạo, thu hồi và tái lập lịch Pod theo yêu cầu tải, khiến địa chỉ IP thay đổi liên tục. Cơ chế kiểm soát truy cập dựa trên định tuyến IP tĩnh truyền thống vì vậy không thể áp dụng trong môi trường này.

    \item \textbf{Lỗ hổng từ các thành phần điều phối:}
    Các giao diện nội bộ của Kubernetes như Kubelet API hoặc cơ chế phân quyền RBAC nếu không được cấu hình chặt chẽ có thể trở thành mục tiêu để kẻ tấn công thao túng cấu hình cụm hoặc leo thang đặc quyền ở tầng điều phối.
\end{itemize}

%------------------------------------------------------------------------------
\subsection{Phân tích bề mặt tấn công}
%------------------------------------------------------------------------------

Để đánh giá toàn diện bề mặt rủi ro từ cả hai tầng trên, đồ án đối chiếu hệ thống với khung \textbf{MITRE ATT\&CK dành cho Container} \parencite{mitre_containers}. Các chiến thuật tấn công được phân loại thành hai giai đoạn nhằm xác định lớp kiểm soát Zero Trust tương ứng:

\begin{table}[H]
\centering
\caption{Giai đoạn 1 -- Xâm nhập ban đầu và thiết lập chỗ đứng}
\label{tab:mitre_phase1}
\renewcommand{\arraystretch}{1.3}
\begin{tabularx}{\textwidth}{|p{3.5cm}|X|p{3.5cm}|}
\hline
\textbf{Giai đoạn} & \textbf{Chiến thuật tiêu biểu} & \textbf{Tầng bị ảnh hưởng} \\
\hline
Xâm nhập ban đầu &
Khai thác lỗ hổng ứng dụng web (T1190); sử dụng tài khoản hợp lệ bị rò rỉ (T1078). &
Tầng ứng dụng (L7) \\
\hline
Thực thi &
Triển khai image chứa mã độc (T1610); thực thi lệnh trái phép qua Kubernetes API (T1609). &
Tầng điều phối \\
\hline
Duy trì truy cập &
Tạo tác vụ ngầm định kỳ (T1053); triển khai container chạy nền (T1525). &
Tầng runtime \\
\hline
Ẩn mình &
Trá hình tiến trình (T1036); xóa nhật ký hệ thống (T1070). &
Tầng runtime \\
\hline
Vô hiệu hóa phòng thủ &
Sửa đổi cấu hình bảo mật đang chạy trong cụm (T1562). &
Tầng runtime \\
\hline
\end{tabularx}
\end{table}

\begin{table}[H]
\centering
\caption{Giai đoạn 2 -- Hậu khai thác và tác động hệ thống}
\label{tab:mitre_phase2}
\renewcommand{\arraystretch}{1.3}
\begin{tabularx}{\textwidth}{|p{3.5cm}|X|p{3.5cm}|}
\hline
\textbf{Giai đoạn} & \textbf{Chiến thuật tiêu biểu} & \textbf{Phạm vi ảnh hưởng} \\
\hline
Leo thang đặc quyền &
Khai thác đặc quyền container để thoát và can thiệp máy chủ vật lý (T1611). &
Tầng nhân hệ điều hành \\
\hline
Đánh cắp thông tin xác thực &
Trích xuất thông tin xác thực tĩnh từ tệp cấu hình hoặc biến môi trường (T1552). &
Tầng lưu trữ thông tin xác thực \\
\hline
Trinh sát &
Dò quét mạng nội bộ để tìm cổng và dịch vụ đang mở (T1046, T1613). &
Tầng mạng (L3/L4) \\
\hline
Di chuyển ngang &
Lợi dụng dịch vụ liên kết nội bộ (T1021) hoặc tái sử dụng token để truy cập chéo. &
Tầng mạng + định danh \\
\hline
Tác động &
Chiếm dụng tài nguyên tính toán (T1496); tấn công từ chối dịch vụ (T1498). &
Tầng runtime + mạng \\
\hline
\end{tabularx}
\end{table}

Phân tích cho thấy các chiến thuật giai đoạn hậu khai thác trải rộng trên nhiều tầng kỹ thuật khác nhau -- từ tầng mạng, tầng runtime đến tầng lưu trữ thông tin xác thực. Không có điểm kiểm soát đơn lẻ nào bao phủ được toàn bộ bề mặt này, đặt ra yêu cầu triển khai các cơ chế phòng thủ đồng thời ở nhiều cấp độ.

%------------------------------------------------------------------------------
\subsection{Các thách thức kỹ thuật}
%------------------------------------------------------------------------------

Từ phân tích trên, việc áp dụng Zero Trust vào môi trường Microservices đặt ra bốn thách thức kỹ thuật cần giải quyết:

\textbf{Định danh độc lập với địa chỉ mạng:}
Do địa chỉ IP thay đổi liên tục, hệ thống cần cơ chế định danh gắn trực tiếp vào bản thân thực thể, dựa trên mật mã học, hỗ trợ xác thực lẫn nhau độc lập với cấu trúc định tuyến. Cơ chế này cần áp dụng cho cả định danh người dùng lẫn định danh workload.

\textbf{Quản lý vòng đời thông tin xác thực ở quy mô lớn:}
Số lượng thông tin xác thực tỷ lệ thuận với số dịch vụ. Hệ thống cần cơ chế cấp phát động, tự động thu hồi và kiểm toán tập trung thay vì lưu trữ tĩnh, để thông tin xác thực không tồn tại lâu hơn thời gian cần thiết.

\textbf{Thực thi chính sách tại nhiều tầng:}
Không một cơ chế kiểm soát đơn lẻ nào bao phủ được toàn bộ bề mặt tấn công. Rủi ro xuất hiện ở tầng giao vận, tầng mạng, tầng giao thức ứng dụng và tầng nhân hệ điều hành. Kiến trúc cần tích hợp chuỗi điểm thực thi chính sách hoạt động đồng thời ở nhiều cấp độ.

\textbf{Ngữ cảnh định danh trong giám sát:}
Trong môi trường IP động, bản ghi lưu lượng mạng thuần túy không đủ để truy vết nguồn gốc luồng dữ liệu. Hệ thống giám sát cần tự động ánh xạ thông tin gói tin với định danh thực thể, namespace và dữ liệu lỗ hổng tương ứng để phục vụ các quyết định an ninh.




%==============================================================================
% 2.2 KHUNG KIẾN TRÚC ZERO TRUST 5 THÀNH PHẦN
%==============================================================================
\section{Đề xuất khung kiến trúc Zero Trust}

Dựa trên phân tích mô hình đe dọa và ba thách thức kỹ thuật cốt lõi đã nêu ở phần trước, đồ án đề xuất một \textbf{Khung kiến trúc Zero Trust} dành riêng cho môi trường Microservices trên Kubernetes. Khung thiết kế này dịch chuyển trọng tâm bảo mật từ ranh giới mạng tĩnh sang việc bảo vệ trực tiếp thực thể và dữ liệu. 

Mọi luồng giao tiếp trong hệ thống không còn dựa vào sự tin cậy ngầm định mà phải tuân thủ nghiêm ngặt quy trình khép kín: \textit{Định danh $\rightarrow$ Đánh giá trạng thái $\rightarrow$ Xác thực $\rightarrow$ Cấp quyền $\rightarrow$ Giám sát và Phản hồi}.

\begin{figure}[H]
\centering
\begin{tikzpicture}[
  font=\footnotesize\sffamily,
  blk/.style={draw, rounded corners=2pt, minimum height=12mm, text width=11.5cm, align=left, inner sep=5pt, line width=0.6pt},
  l1/.style={blk, fill=blue!8, draw=blue!50},
  l2/.style={blk, fill=orange!12, draw=orange!60},
  l3/.style={blk, fill=green!10, draw=green!50},
  l4/.style={blk, fill=red!10, draw=red!50},
  l5/.style={blk, fill=violet!10, draw=violet!60}
  ]
\node[l1] (l1) at (0, 10) {
  \textbf{Thành phần 1 — Định danh đa thực thể:}\\[3pt]
  Xác thực người dùng cuối
  | Cấp định danh cho khối lượng công việc (Workload) 
  | Gốc tin cậy hạ tầng};

\node[l2] (l2) at (0, 7.5) {
  \textbf{Thành phần 2 — Đánh giá tư thế và tình báo đe dọa:}\\[3pt]
  Kiểm toán lỗ hổng mã nguồn 
  | Kiểm soát cấu hình tại cổng nạp 
  | Cập nhật tri thức mối đe dọa bên ngoài};

\node[l3] (l3) at (0, 5) {
  \textbf{Thành phần 3 — Thực thi chính sách đa tầng:}\\[3pt]
  Kiểm soát truy cập biên 
  | Phân đoạn vi mô nội bộ
  | Giám sát hệ thống tại thời điểm chạy (Runtime)};

\node[l4] (l4) at (0, 2.5) {
  \textbf{Thành phần 4 — Quản lý thông tin xác thực động:}\\[3pt]
  Cấp phát thông tin xác thực theo cơ chế dùng một lần (JIT) 
  | Tự động thu hồi};

\node[l5] (l5) at (0, 0) {
  \textbf{Thành phần 5 — Quan sát và phản hồi thích ứng:}\\[3pt]
  Thu thập nhật ký tập trung 
  | Ánh xạ ngữ cảnh mạng 
  | Đánh giá và cập nhật điểm tin cậy liên tục};

\draw[->, >=Latex, dashed, gray!80, line width=0.8pt] (l5.east) -- ++(1.2,0) |- node[pos=0.25, right, font=\scriptsize\sffamily, align=left, xshift=2pt] {Vòng phản hồi\\(Cập nhật điểm tin cậy\\và điều chỉnh quyền truy cập)} (l1.east);
\end{tikzpicture}
\caption{Khung kiến trúc logic Zero Trust 5 thành phần đề xuất}
\label{fig:logical_5layer_design}
\end{figure}

Khung kiến trúc được thiết kế nhằm giải quyết trực tiếp các thách thức kỹ thuật đã đặt ra. Dưới đây là chức năng và vai trò của từng thành phần:

\subsection{Thành phần 1: Định danh đa thực thể}
Thành phần này đóng vai trò là gốc tin cậy (Root of Trust) cho toàn bộ hệ thống, giải quyết trực tiếp thách thức về sự vô nghĩa của địa chỉ IP động. Kiến trúc áp dụng cơ chế \textbf{định danh kép}:
\begin{itemize}[itemsep=6pt, parsep=0pt]
    \item \textbf{Định danh người dùng:} Xác thực người dùng tại biên và đóng gói quyền hạn vào các mã thông báo bảo mật để phục vụ ủy quyền ở các tầng sau.
    \item \textbf{Định danh thực thể (Workload):} Mọi vi dịch vụ khi khởi tạo đều được cấp một danh tính mật mã học duy nhất có thời gian sống rất ngắn. Thực thể sử dụng danh tính này để xác thực chéo lẫn nhau (mTLS), loại bỏ hoàn toàn sự phụ thuộc vào mạng định tuyến định tĩnh.
\end{itemize}

\subsection{Thành phần 2: Đánh giá tư thế bảo mật và tình báo đe dọa}
Chỉ có định danh hợp lệ là chưa đủ, hệ thống cần đánh giá mức độ an toàn của hệ thống trước khi quyết định cấp quyền. Thành phần này hoạt động như một Điểm cung cấp thông tin (PIP), liên tục thu thập ngữ cảnh bảo mật cho Bộ xử lý chính sách:
\begin{itemize}[itemsep=6pt, parsep=0pt]
    \item \textbf{Tư thế của Workload:} Liên tục quét mã nguồn và ảnh vùng chứa (container image) để tìm kiếm các lỗ hổng bảo mật hoặc cấu hình sai lệch.
    \item \textbf{Tình báo đe dọa:} Tự động đồng bộ danh sách các địa chỉ IP và tên miền độc hại từ Internet, giúp hệ thống phòng vệ chủ động ngăn chặn các nỗ lực giao tiếp với máy chủ điều khiển (Command \& Control).
\end{itemize}

\subsection{Thành phần 3: Thực thi chính sách đa tầng}
Để khắc phục rủi ro không có điểm thực thi chính sách đơn nhất, thành phần này thiết lập chiến lược phòng thủ theo chiều sâu với các Điểm thực thi chính sách (PEP) hoạt động ở 4 cấp độ độc lập, dựa trên nguyên tắc kiểm soát truy cập phân quyền theo thuộc tính (ABAC):
\begin{itemize}[itemsep=6pt, parsep=0pt]
    \item \textbf{Nguyên tắc từ chối mặc định (Default Deny):} Mọi luồng giao tiếp mặc định bị chặn. Hệ thống chỉ cho phép các kết nối được khai báo tường minh.
    \item \textbf{Tầng mạng biên (North-South):} Đứng tại cửa ngõ hệ thống, kiểm tra chữ ký và hiệu lực của mã thông báo truy cập.
    \item \textbf{Tầng mạng nội bộ (East-West):} Thực thi chính sách phân đoạn vi mô (microsegmentation) tại tầng mạng và giao thức ứng dụng. Mỗi luồng giao tiếp giữa các dịch vụ chỉ được thiết lập khi xác thực thành công định danh mật mã học của thực thể nguồn.
    \item \textbf{Tầng cấu hình và Runtime:} Chặn việc triển khai các ứng dụng không có chữ ký số hợp lệ và ngắt ngay lập tức các tiến trình thực thi lời gọi hệ thống (syscall) bất thường từ bên trong vùng chứa.
\end{itemize}

\subsection{Thành phần 4: Quản lý thông tin xác thực động}
Thành phần này giải quyết thách thức phân mảnh thông tin xác thực tĩnh. Kiến trúc loại bỏ việc lưu trữ mật khẩu tĩnh trong tệp cấu hình, thay vào đó áp dụng cơ chế cấp phát tức thời (Just-in-Time). Khi một vi dịch vụ cần truy cập cơ sở dữ liệu, Bộ điều phối chính sách sẽ cấp một tài khoản ngẫu nhiên có thời gian sống (TTL) ngắn. Khi hết hạn, mật khẩu tự động bị thu hồi. Các bí mật này chỉ được lưu trữ trên bộ nhớ ảo của vùng chứa thay vì ổ cứng vật lý.

\subsection{Thành phần 5: Quan sát và phản hồi thích ứng}
Hệ thống giám sát trong Zero Trust không chỉ lưu trữ nhật ký để hậu kiểm mà đóng vai trò cung cấp dữ liệu đầu vào trực tiếp để ra quyết định \parencite{nist800207}. Khung kiến trúc ứng dụng một mô hình đánh giá điểm tin cậy dựa trên rủi ro (Risk-based Trust Algorithm):
\begin{itemize}[itemsep=6pt, parsep=0pt]
    \item \textbf{Quan sát có ngữ cảnh:} Thu thập toàn bộ luồng mạng và sự kiện bảo mật. Hệ thống tự động ánh xạ thông tin gói tin mạng với định danh thực thể và không gian tên tương ứng.
    \item \textbf{Vòng lặp phản hồi thích ứng:} Kiến trúc áp dụng thuật toán tính điểm tin cậy dựa trên các yếu tố phạt (penalty-based). Khi phát hiện một vùng chứa xuất hiện lỗ hổng nghiêm trọng hoặc có hành vi gọi mạng bất thường, hệ thống tự động đánh tụt điểm tin cậy của thực thể đó. Sự thay đổi trạng thái này lập tức kích hoạt các PEP ở Thành phần 3 để cắt đứt quyền truy cập và cô lập mối đe dọa, loại bỏ hoàn toàn độ trễ của việc can thiệp thủ công.
\end{itemize}







































%==============================================================================
% 2.2 KHUNG KIẾN TRÚC ZERO TRUST 5 THÀNH PHẦN
%==============================================================================
\section{Đề xuất khung kiến trúc Zero Trust 5 thành phần}

Dựa trên phân tích mô hình đe dọa và các thách thức kỹ thuật, đồ án đề xuất một \textbf{Khung kiến trúc Zero Trust 5 thành phần} dành riêng cho môi trường Microservices trên Kubernetes. 

Khung kiến trúc đề xuất hiện thực hóa ba vai trò logic cốt lõi của tiêu chuẩn NIST SP 800-207 (PIP, PDP, PEP) thành năm thành phần chức năng, được phân tách dựa trên cơ sở tần suất cập nhật dữ liệu và vị trí của chúng trong luồng xử lý \parencite{nist800207}. Việc phân rã mô hình lý thuyết này giúp kiến trúc có thể triển khai thực tế, cụ thể:

\begin{itemize}[itemsep=6pt, parsep=0pt]
    \item \textbf{Vai trò Điểm cung cấp thông tin (PIP)} có phạm vi quá rộng trong lý thuyết, do đó được tách thành ba nhóm dữ liệu độc lập: \textit{Thành phần 1} quản lý định danh (tần suất cập nhật thấp), \textit{Thành phần 2} quản lý trạng thái lỗ hổng (tần suất cập nhật trung bình), và \textit{Thành phần 5} thu thập nhật ký mạng (thời gian thực).
    
    \item \textbf{Vai trò Điểm quyết định chính sách (PDP)} được cấu thành từ \textit{Thành phần 4} (đóng vai trò Quản trị chính sách - PA trong việc phân phối bí mật động) và lõi phân tích của \textit{Thành phần 5} (đóng vai trò Động cơ chính sách - PE để tính toán rủi ro).
    
    \item \textbf{Vai trò Điểm thực thi chính sách (PEP)} được đối chiếu trực tiếp thành \textit{Thành phần 3}, nhưng được mở rộng thành một chuỗi kiểm soát đa tầng thay vì một điểm kiểm soát duy nhất.
\end{itemize}

Mọi luồng giao tiếp trong hệ thống phải tuân thủ quy trình khép kín: \textit{Định danh $\rightarrow$ Đánh giá rủi ro $\rightarrow$ Xác thực $\rightarrow$ Cấp quyền $\rightarrow$ Giám sát và Phản hồi}.

\begin{figure}[H]
\centering
\begin{tikzpicture}[
  font=\footnotesize\sffamily,
  blk/.style={draw, rounded corners=2pt, minimum height=12mm, text width=11.5cm, align=left, inner sep=5pt, line width=0.6pt},
  l1/.style={blk, fill=blue!8, draw=blue!50},
  l2/.style={blk, fill=orange!12, draw=orange!60},
  l3/.style={blk, fill=green!10, draw=green!50},
  l4/.style={blk, fill=red!10, draw=red!50},
  l5/.style={blk, fill=violet!10, draw=violet!60}
  ]
\node[l1] (l1) at (0, 10) {
  \textbf{Thành phần 1 — Định danh đa thực thể (Ngữ cảnh PIP):}\\[3pt]
  Xác thực người dùng cuối
  | Cấp định danh mật mã học cho khối lượng công việc};

\node[l2] (l2) at (0, 7.5) {
  \textbf{Thành phần 2 — Đánh giá tư thế và tình báo đe dọa (Dữ liệu PIP):}\\[3pt]
  Quét lỗ hổng mã nguồn 
  | Kiểm soát cấu hình 
  | Cập nhật tri thức mối đe dọa};

\node[l3] (l3) at (0, 5) {
  \textbf{Thành phần 3 — Thực thi chính sách đa tầng (PEP):}\\[3pt]
  Kiểm soát truy cập biên 
  | Phân đoạn vi mô nội bộ
  | Giám sát hệ thống thời điểm chạy};

\node[l4] (l4) at (0, 2.5) {
  \textbf{Thành phần 4 — Quản lý thông tin xác thực động (Hỗ trợ PA):}\\[3pt]
  Cấp phát thông tin xác thực theo cơ chế dùng một lần (JIT) 
  | Tự động thu hồi};

\node[l5] (l5) at (0, 0) {
  \textbf{Thành phần 5 — Đánh giá rủi ro và Phản hồi thích ứng (Trung tâm PE \& PIP):}\\[3pt]
  Thu thập nhật ký 
  | Đánh giá mức độ nghiêm trọng 
  | Điều chỉnh chính sách bảo mật};

\draw[->, >=Latex, dashed, gray!80, line width=0.8pt] (l5.east) -- ++(1.2,0) |- node[pos=0.25, right, font=\scriptsize\sffamily, align=left, xshift=2pt] {Vòng phản hồi\\(Cập nhật quyết định\\và điều chỉnh truy cập)} (l1.east);
\end{tikzpicture}
\caption{Khung kiến trúc Zero Trust 5 thành phần và đối chiếu vai trò logic NIST}
\label{fig:logical_5layer_design}
\end{figure}